\begin{textsample}{NEG Dim 2 – 2022\_c5 – Score -1.00 – 2022\_c5/enade\_2022\_tec\_mark\_64d0806439.txt}  \label{ex:f2_neg_217}
CAPÍTULO 6 DISTRIBUIÇÃO DE CONCEITOS A finalidade deste capítulo é apresentar um panorama nacional sobre a distribuição do Conceito Enade 2022.
Para traçar esse panorama, recorre -se a tabelas e gráficos, que são acompanhados de comentários e análises que contribuem para a contextualização dos dados neles expostos.
Para visualizar o panorama do Enade/2002, em nível nacional, na primeira seção do capítulo, os dados sobre o número de cursos participantes são organizados por grande região.
Na segunda, o mapeamento do conceito do exame é feito por categoria administrativa das IES e modalidade de oferta dos cursos, de maneira cruzada às grandes regiões.
Na última seção, a lógica de apresentação dos dados é feita por organização acadêmica e, também, grande região de funcionamento dos cursos, no país como um todo. 6.1 PANORAMA NACIONAL DA DISTRIBUIÇÃO DOS CONCEITOS A lógica adotada na apresentação dos dados, nesta seção e nas duas seguintes, pressupõe -se o registro de duas informações de natureza reiterativa.
A primeira informação diz respeito aos valores possíveis do Enade, que variam em níveis de 1 a 5, sendo que, à medida que esses valores aumentam, o curso recebe uma avaliação melhor ; além disso, os cursos com somente um ou sem nenhum concluinte participante não obtêm o Conceito Enade, ficando Sem Conceito ( SC ), como ressaltado na Apresentação deste relatório.
O segundo registro é que a diferença entre os cursos tabulados neste capítulo e no Capítulo 3 corresponde aos cursos que não foram avaliados, ou seja, àqueles, em princípio, sem estudantes concluintes que tivessem sido inscritos para a prova.
Na Tabela 6.1, são apresentadas a quantidade e a distribuição de cursos de Tecnologia em Marketing participantes do Enade/2022, por faixa de conceito e grande região.
Observando -se os dados da Tabela 6.1, nota -se que, dos 344 cursos participantes, 140 ( 40,7 \% ) classificaram -se com conceito 3, o valor modal.
Esse foi também o conceito modal em todas as regiões : Nordeste ( 41,8 \% ), Sudeste ( 37,1 \% ), Sul ( 46,7 \% ), Norte ( 58,8 \% ) e Centro-Oeste ( 38,9 \% ).
O conceito 2 foi o segundo mais frequente em nível nacional ( 16,9 \%, correspondendo a 58 cursos ).
O conceito 4 foi o terceiro ( 13,7 \%, correspondendo a 47 cursos ).
Houve, ainda, 13 cursos ( 3,8 \% ) que receberam o conceito 5, e 7 cursos ( 2,0 \% ) que receberam o conceito 1.
Dos 344 cursos de Tecnologia em Marketing, 79 ( 23 \% ) ficaram Sem Conceito ( SC ).
A região Norte participou com 17 cursos, o que corresponde a 4,9 \% do total nacional.
Desses, 2 cursos receberam o conceito 2, o que equivale a 11,8 \% do total regional, o conceito modal na região, como já comentado.
A nenhum curso atribuiu -se o conceito 1, e a 10 cursos ( 58,8 \% ), o conceito 3.
Quatro cursos ( 23,5 \% ) ficaram Sem Conceito ( SC ), um curso ( 5,9 \% ) recebeu o conceito 5, e nenhum curso o conceito 4, nesta região.
A região Nordeste participou com 55 cursos, o correspondente a 15,9 \% do total nacional.
Nessa região, 23 cursos ( 41,8 \% em termos regionais ) obtiveram conceito 3, o conceito modal para a região, como já comentado.
Os conceitos 1 e 2 foram atribuídos, respectivamente, a 1 e a 7 cursos ( 1,8 \% e 12,7 \%, respectivamente ).
Ainda nessa região, 7 cursos ( 12,7 \% ) foram avaliados com o conceito 4, e nenhum curso foi avaliado com o conceito 5.
Ressaltamos que 17 cursos ( 30,9 \% ) ficaram Sem Conceito ( SC ).
Para os 194 cursos participantes da região Sudeste ( 56,4 \% do total nacional ), o conceito 3 foi o conceito modal, como já anteriormente assinalado, tendo sido obtido por 72 cursos ( 37,1 \% ).
O conceito 1 foi atribuído a 5 cursos ( 2,6 \% ), e o conceito 2, a 41 cursos ( 21,1 \% ).
Receberam os conceitos 4 e 5, respectivamente, 28 cursos ( 14,4 \% ) e 6 cursos ( 3,1 \% ).
Nessa região, 42 cursos ficaram Sem Conceito ( SC ).
Os 60 cursos da região Sul corresponderam a 17,4 \% do total nacional.
Como já indicado, houve predominância do conceito 3 ( conceito modal ), atribuído a 46,7 \% dos cursos da região, ou seja, atribuído a 28 dos 60 cursos participantes na região Sul.
O conceito 1 não foi atribuído a nenhum curso ( 0,0 \% ) ; já o conceito 2, a 7 cursos ( 11,7 \% ) ; e os conceitos 4 e 5, a 9 ( 15 \% ) e a 6 cursos ( 10 \% ), respectivamente.
Na região Sul, dez cursos ficaram Sem Conceito ( SC ).
Já dos 18 cursos participantes na região Centro-Oeste ( 5,2 \% do total nacional ), 7 ( 38,9 \% em termos regionais ) receberam conceito 3, o conceito modal, como já destacado.
Um curso ( 5,6 \% ) obteve o conceito 1 ; um curso ( 5,6 \% ) obteve o conceito 2 ; 3 ( 16,7 \% ), o conceito 4 ; e nenhum, o conceito 5.
Nessa região, 6 cursos ficaram Sem Conceito ( SC ).
Em um gráfico de distribuição cumulativa, em cada valor de abscissa ( eixo x ), o valor apresentado no eixo das ordenadas ( eixo y ) é o valor acumulado da variável ( Conceito Enade ).
Em linhas gerais, em um gráfico desta mesma tipologia, poligonais mais à esquerda representam, no caso dos dados apresentados nesta seção, as grandes regiões com a pior distribuição, e poligonais mais à direita, as grandes regiões com os melhores conceitos.
A partir destas diretrizes, no Gráfico 6.1, é apresentada a distribuição cumulativa do Conceito Enade, segundo a grande região.
Nesse gráfico, poligonais mais à esquerda representam Conceito Enade mais baixos, e poligonais mais à direita, Conceito Enade mais altos. 6.2 CONCEITOS POR CATEGORIA ADMINISTRATIVA, MODALIDADE DE OFERTA DOS CURSOS E GRANDE REGIÃO Na Tabela 6.2, é apresentada a distribuição dos cursos de Tecnologia em Marketing participantes do Enade/2022, por categoria administrativa e modalidade de oferta, de acordo com os conceitos por eles alcançados, segundo as grandes regiões brasileiras.
De acordo com as informações da Tabela 6.2, dos 344 cursos participantes, 9 ( 2,6 \% ) eram ministrados em instituições públicas, e 335 ( 97,3 \% ), em instituições privadas.
Quanto à modalidade de oferta, 168 cursos eram a distância, e 176 eram presenciais.
Ainda conforme as informações da mesma tabela, em termos nacionais, entre instituições públicas e privadas, observa -se que, dos 13 cursos avaliados com conceito 5, 1 foi oferecido em IES públicas, e 12, em IES privadas.
Dos 9 cursos participantes de IES públicas, o conceito 4 foi o valor modal, atribuído a 6 cursos ( 66,6 \% ).
Dos demais cursos de instituições públicas participantes, 2 cursos ( 22,2 \% ) foram avaliados com o conceito 3, e 1 curso ( 11,1 \% ) foi avaliado com o conceito 5.
Nessa categoria, nenhum curso ficou Sem Conceito ( SC ) ou recebeu conceito 1 e 2.
Na rede privada, o conceito modal foi o 3, atribuído a 138 cursos, o correspondente a 41,1 \% dos 335 cursos da categoria.
Dos demais cursos participantes, 7 ( 2 \% ) receberam o conceito 1, e 58 ( 17,3 \% ), o conceito 2.
O conceito 4 foi atribuído a 41 cursos ( 12,2 \% ), e o conceito 5, a 21 cursos ( 2,3 \% ).
Nessa categoria administrativa, 79 cursos ( 23,5 \% ) ficaram Sem Conceito ( SC ).
Dos 176 cursos participantes da modalidade presencial, o conceito 3 foi o valor modal, atribuído a 58 cursos ( 32,9 \% ).
Dos demais cursos nessa modalidade de oferta, 3 cursos ( 1,7 \% ) receberam o conceito 1 ; 29 cursos ( 16,4 \% ) receberam o conceito 2 ; 29 cursos ( 16,4 \% ), o conceito 4 ; e 10 cursos ( 5,6 \% ), o conceito 5.
Nessa modalidade de oferta, 47 cursos ( 26,7 \% ) ficaram Sem Conceito ( SC ).
Dos 168 cursos participantes da modalidade a distância, o conceito 3 foi o valor modal, atribuído a 82 cursos ( 48,8 \% ).
Dos demais cursos nessa modalidade de oferta, 4 cursos ( 2,3 \% ) receberam o conceito 1 ; 29 cursos ( 17,2 \% ) receberam o conceito 2 ; 18 cursos, ( 10,7 \% ), o conceito 4 ; e 3 cursos ( 1,7 \% ), o conceito 5.
Nessa modalidade de oferta, 32 cursos ( 19 \% ) ficaram Sem Conceito ( SC ).
Na análise por região, observa -se que, na região Norte, as instituições privadas participaram com 17 cursos ( 4,9 \% do total regional ), e 4 cursos ficaram Sem Conceito ( SC ).
O conceito modal para as instituições privadas na região foi o 3, atribuído a 10 cursos, correspondendo a 58,8 \% dessa categoria na região.
Nenhum curso foi avaliado com o conceito 1 ou 4 ; 2, com o conceito 2 ; e um, com conceito 5.
As instituições públicas não ofereceram nenhum curso na região Norte.
Dos cursos oferecidos na região Norte, 9 eram presenciais, e 3 ficaram Sem Conceito ( SC ).
Dos 9 cursos oferecidos, nenhum obteve o conceito 1 ou 4 ; 1 curso, o conceito 2 ; 10 cursos obtiveram o conceito 3 ; e 1, o conceito 5. 8 eram na modalidade a distância, tendo 6 recebido o conceito modal 3 ; 1, o conceito 2 ; e 1 ficou Sem Conceito ( SC ).
Nenhum curso recebeu os demais conceitos.
Na região Nordeste, a rede privada concentrou 54 dos 55 cursos participantes, o equivalente a 98,1 \% do total da região.
Entre os cursos oferecidos por IES privadas no Nordeste, 22 cursos foram avaliados com conceito 3, o conceito modal. 1 curso obteve o conceito 1 ; 7, o conceito 2 ; 7, o conceito 4 ; e nenhum curso obteve o conceito 5.
Nessa categoria administrativa e grande região, 17 cursos ficaram Sem Conceito ( SC ).
As instituições públicas da região Nordeste participaram com 1 curso ( 1,8 \% do total da região ), que recebeu o conceito 3.
Dos cursos oferecidos na região Nordeste, 28 eram presenciais : 3 ficaram Sem Conceito ( SC ), um ficou com o conceito 1 ; 5, com o conceito 2 ; 10, com o conceito 3, o valor modal ; 3, com o conceito 4 ; e nenhum com o conceito 5.
Na modalidade a distância, eram 27 cursos, tendo nenhum obtido o conceito 1 ; 2, o conceito 2 ; 13, o conceito 3, o valor modal ; 4, o conceito 4.
Nessa região, nenhum curso ficou Sem Conceito ( SC ) ou recebeu o conceito 5.
Na região Sudeste, a proporção de cursos da rede privada, 97,4 \%, correspondeu a 189 dos 194 cursos participantes.
Nessa combinação de categoria administrativa e grande região, o conceito modal foi o 3, atribuído a 71 cursos. 42 cursos ficaram Sem Conceito ( SC ).
Os demais foram avaliados com os seguintes conceitos : 1 ( 5 cursos ), 2 ( 41 cursos ), 4 ( 25 cursos ) e 5 ( 5 cursos ).
Dos 5 cursos oferecidos por instituições públicas na região Sudeste, o conceito modal foi o 4, atribuído a 3 cursos.
Dos demais cursos, 1 recebeu o conceito 3, e 1 recebeu o conceito 5.
Nessa combinação de categoria administrativa e grande região, nenhum curso ficou Sem Conceito ( SC ), e nenhum recebeu os conceitos 1 e 2.
Também na região Sudeste, grande parcela dos cursos foi oferecida na modalidade presencial : 111 de 194 cursos.
Para esses, a moda foi o conceito 3, com 35 cursos.
Os demais cursos presenciais da região foram alocados nos conceitos 1 ( 2 cursos ), 2 ( 23 cursos ), 4 ( 19 cursos ) e 5 ( 4 cursos ).
Ficaram Sem Conceito ( SC ) 28 cursos presenciais na região Sudeste.
Na região Sudeste, 83 cursos foram oferecidos na modalidade a distância, tendo sido 3 cursos com conceito 1 ; 18 com conceito 2 ; 37 com conceito 3, o valor modal ; 9 com o conceito 4 ; e 2 com o conceito 5.
Na região Sudeste, 14 cursos a distância ficaram Sem Conceito ( SC ).
As instituições privadas concentraram 58 dos 60 cursos participantes da região Sul, 96,6 \% do total regional.
Desses, 28 cursos obtiveram conceito 3, o conceito modal.
Nessa combinação de categoria administrativa e grande região, nenhum curso recebeu o conceito 1 ; 7 receberam o conceito 2 ; outros 7, o conceito 4 ; e ainda 6, o conceito 5.
Dez cursos ficaram Sem Conceito ( SC ).
As instituições públicas na região Sul participaram com 2 cursos ( 3,4 \% dos cursos da região ), aos quais foram atribuídos apenas o conceito 4 ( 2 cursos ).
Nenhum curso da região Sul oferecido por IES públicas ficou Sem Conceito ( SC ) ou recebeu os demais conceitos.
Na região Sul, 23 dos cursos foram oferecidos na modalidade presencial.
Para esses, a moda foi o conceito 3, com 7 cursos.
Os demais cursos presenciais da região foram alocados nos conceitos 4 ( 6 cursos ) e 5 ( 5 cursos ).
Ficaram Sem Conceito ( SC ) cinco cursos e nenhum curso recebeu os conceitos 1 e 2.
Na modalidade a distância, foram oferecidos 37 cursos, sendo 7 cursos com conceito 7 ; 21 com conceito 3, o valor modal ; 3 com conceito 4 ; e um com o conceito 5.
Nenhum curso ficou Sem Conceito ( SC ) ou com o conceito 1 nesta região.
Na região Centro-Oeste, 17 dos 18 cursos participantes eram de IES privadas ( 94,4 \% em termos regionais ).
Desses, 7 cursos receberam conceito 3, o valor modal.
Dos demais cursos, um recebeu o conceito 1 ; um recebeu o conceito 2 ; 2, o conceito 4 ; e nenhum curso recebeu o conceito 5.
Seis cursos ficaram Sem Conceito ( SC ).
Quanto ao único curso oferecido por instituições públicas na região Centro-oeste ( 5,6 \% do total regional ), apresentou o conceito 4.
Também na região Centro-oeste, a maioria dos cursos foram a distância : 13 ( 72,2 \% ).
Para esses, a moda foi o conceito 3, com 5 cursos ( 38,4 \% ).
Os demais cursos a distância da região foram alocados nos conceitos 1 ( um curso ), 2 ( um curso ) e 4 ( 2 cursos ).
Quatro cursos a distância da região ficaram Sem Conceito ( SC ).
Para os 5 cursos presenciais ( 27,7 \% ), a moda unimodal foi o conceito 3 ( 2 cursos ).
Os demais cursos a distância da região receberam os conceitos 4 ( um curso ).
Dois cursos presenciais da região ficaram Sem Conceito ( SC ) e nenhum ficou com o conceito 1, 2 e 5.
A título de síntese, no Gráfico 6.2, apresenta -se a distribuição do Conceito Enade, segundo a categoria administrativa da IES.
Nesse gráfico, como dito anteriormente, poligonais mais à esquerda representam Conceito Enade mais baixos, e poligonais mais à direita, valores mais altos.
É importante observar que, para cada poligonal presente nesse gráfico, a porcentagem constante no eixo y se refere ao universo de cursos oferecidos por instituições públicas ( poligonal vermelha ) ou privadas ( poligonal azul ). 6.3 CONCEITOS POR ORGANIZAÇÃO ACADÊMICA E GRANDE REGIÃO Na Tabela 6.3, encontra -se a distribuição dos conceitos atribuídos aos cursos participantes do Enade/2022, na área de Tecnologia em Marketing, por organização acadêmica, segundo as grandes regiões brasileiras.
Como mostram os dados da Tabela 6.3, dos 344 cursos de Tecnologia em Marketing participantes, 121 eram oferecidos em universidades ; 162, em centros universitários ; 58 em faculdades ; e 3, em CEFET/IF.
Esta distribuição corresponde a, respectivamente, 35,1 \%, 48,8 \%, 16,8 \% e 0,8 \% dos cursos.
De acordo com os dados apresentados na mesma tabela, de todos os 13 cursos avaliados com o conceito 5, 2 eram vinculados a universidades.
Os cursos vinculados a esse tipo de organização acadêmica tiveram o conceito 3 como conceito modal, atribuído a 48 cursos ( 39,6 \% ).
Os demais cursos ofertados por universidades receberam os conceitos 1 ( um curso ), 2 ( 28 cursos ), 4 ( 17 cursos ) e conceito 5 ( 2 cursos, como já mencionado ).
Vinte e cinco cursos ficaram Sem Conceito ( SC ).
Dos cursos vinculados a centros universitários, o conceito modal também foi o 3, atribuído a 72 cursos ( 44,4 \% ). 34 cursos vinculados a esse tipo de organização acadêmica ficaram Sem Conceito ( SC ), e os outros cursos receberam os conceitos 1 ( 5 cursos ), 2 ( 23 cursos ), 4 ( 22 cursos ) e 5 ( 6 cursos ).
Dos 58 cursos mantidos por faculdades, 19 ( 32,7 \% ) receberam conceito 3, o conceito modal.
O segundo conceito mais frequentemente atribuído aos cursos oferecidos por essa categoria de IES foi o conceito 2, recebido por 7 cursos.
Na sequência, 6 cursos obtiveram conceito 4 ; 1 curso, o conceito 1 ; e 5 cursos, o conceito 5.
Dos cursos oferecidos em faculdades, 20 ficaram Sem Conceito ( SC ).
Dos 3 cursos mantidos por CEFET/IF, dois ( 66,6 \% ) receberam conceito 4, o conceito modal.
O segundo conceito mais frequentemente atribuído aos cursos oferecidos por essa categoria de IES foi o conceito 3, recebido por um curso.
Nenhum curso ficou Sem Conceito ( SC ) ou com os conceitos 1, 2 e 5.
Considerando -se, separadamente, as regiões brasileiras, verifica -se que, na região Norte, as universidades participaram com 5 cursos avaliados, aos quais foram atribuídos os seguintes conceitos : 2, a 1 curso ; 3, o conceito modal, a 2 cursos ; e 1, 4 e 5 a nenhum curso.
Dois cursos ficaram Sem Conceito ( SC ) nessa combinação de categoria e região.
Os centros universitários da região Norte foram representados por 7 cursos, dos quais 5 obtiveram conceito 3, o valor modal, e um curso, o conceito 2.
Os conceitos 1, 4 e 5 não foram atribuídos a nenhum curso.
Um curso ficou Sem Conceito ( SC ).
As faculdades participaram com 5 cursos na região Norte, os quais receberam os conceitos 3 ( três cursos ) e 5 ( um curso ).
Os conceitos 1, 2 e 4 não foram atribuídos a nenhum curso.
Um curso ficou Sem Conceito ( SC ).
Nessa região e organização acadêmica, nenhum curso foi oferecido em CEFET/IF.
Na região Nordeste, as universidades participaram com 13 dos 55 cursos da Área de Tecnologia em Marketing oferecidos.
O conceito bimodal foi 3 e 4, atribuído a 4 cursos, cada.
Outros 3 cursos obtiveram o conceito 2 ; e nenhum obteve o conceito 1 e 5.
Dois cursos nas universidades dessa região ficaram Sem Conceito ( SC ).
Os centros universitários contaram com 29 cursos participantes na região Nordeste, 14 dos quais receberam o conceito modal 3.
Para os demais cursos, foram atribuídos os conceitos 2 ( 3 cursos ) e 4 ( 3 cursos ).
Nove cursos oferecidos por centros universitários no Nordeste ficaram Sem Conceito ( SC ) ; nenhum recebeu o conceito 1 e 5.
As faculdades foram representadas por 12 cursos na região Nordeste, 4 dos quais receberam conceito 3, o conceito modal.
Aos demais cursos foram atribuídos os conceitos 1 ( um curso ) e 2 ( um curso ), e nenhum curso recebeu o conceito 4 e 5.
Seis cursos ficaram Sem Conceito ( SC ).
Nessa combinação de região e organização acadêmica, um curso foi oferecido em CEFET/IF, que recebeu o conceito 3.
Na região Sudeste, as universidades concentraram 71 dos 194 cursos de Tecnologia em Marketing da região.
Dos cursos oferecidos em IES com este tipo de organização acadêmica, nessa região, o conceito modal foi 3, atribuído a 27 cursos.
Os demais cursos receberam os conceitos 2 ( 21 cursos ), 4 ( 6 cursos ) e 1 ( um curso ).
Dezesseis cursos ficaram Sem Conceito ( SC ), e nenhum recebeu o conceito 5.
Os centros universitários participaram com 87 cursos na região Sudeste, dos quais 35 obtiveram o conceito modal 3.
O conceito 1 foi atribuído a 4 cursos ; o conceito 2, a 16 cursos ; o conceito 4, a 16 cursos ; e o conceito 5, a 3 cursos.
Quinze cursos ficaram Sem Conceito ( SC ) nessa combinação de organização acadêmica e região.
Já as faculdades foram representadas por 35 cursos na região Sudeste, dos quais onze ficaram Sem Conceito ( SC ).
O conceito modal 3 foi atribuído a 10 cursos.
Aos demais cursos foram atribuídos os conceitos 2 ( 6 cursos ), 4 ( 5 cursos ) e 5 ( 3 cursos ).
Nenhum curso dessa organização acadêmica recebeu o conceito 1.
Para o único curso participante de CEFET/IF, na região Sudeste, o conceito foi o 4.
Dos 60 cursos da região Sul, 25 eram mantidos por universidades, para os quais o conceito modal foi 3, atribuído a 13 cursos.
Os demais cursos receberam os conceitos 2 ( 2 cursos ), 4 ( 5 cursos ), 5 ( 2 cursos ) ; nenhum curso recebeu o conceito 1.
Em IES com esse tipo de organização acadêmica, três cursos da região Sul ficaram Sem Conceito ( SC ).
Dos 29 cursos participantes de centros universitários da região Sul, 5 ficaram Sem Conceito ( SC ), e 14 cursos receberam conceito 3, o valor modal.
Os conceitos restantes ficaram assim \textbf{distribuídos} : nenhum curso ( conceito 1 ) ; 5 cursos ( conceito 2 ) ; 2 cursos ( conceito 4 ) ; e 3 cursos ( conceito 5 ).
Quanto aos 5 cursos vinculados a faculdades na região Sul, 1 curso recebeu o conceito 3, 1 curso o conceito 4 e 1 curso recebeu o conceito 5.
Somente 2 cursos ficaram Sem Conceito ( SC ) e nenhum curso recebeu o conceito 1 e 2.
Do único curso oferecido em CEFET/IF, foi avaliado com o conceito 4.
Na região Centro-Oeste, 7 dos 18 cursos oferecidos eram mantidos por universidades.
Para esse tipo de organização acadêmica, o conceito bimodal foi 3 e 4, atribuído a 2 cursos, cada ; e, dois cursos da região ficaram Sem Conceito ( SC ).
O curso restante recebeu os conceitos 2, os demais conceitos não foram atribuídos a nenhum curso.
Os centros universitários da região Centro-Oeste contaram com 10 cursos : 1 curso foi avaliado com o conceito 1 ; nenhum curso foi avaliado com o conceito 2 ; 4 cursos, com o conceito 3 ( modal ) ; 1 curso recebeu o conceito 4.
Não houve cursos avaliados com o conceito 5, e 2 cursos ficaram Sem Conceito ( SC ).
Do único curso oferecido em faculdades na região Centro-Oeste, esse recebeu o conceito 3.
Não houve cursos oferecidos em CEFET/IF nessa região.
De maneira semelhante à forma como se procedeu, em relação à seção anterior, encerra -se com o Gráfico 6.3, que apresenta a distribuição do Conceito Enade dos cursos de Tecnologia em Marketing segundo a organização acadêmica da IES.
A leitura desse último gráfico é semelhante aos anteriores, seguindo -se a lógica de que poligonais deslocadas para a esquerda correspondem a Conceito Enade mais baixos ( 1, 2 e 3 ) e poligonais deslocadas para a direita correspondem a Conceito Enade mais altos.

% matched lemmas: distribuir
\end{textsample}
